% GENERATED FILE. Edit canonical lessons, then run npm run generate:books.
% canonical-source-hash: fnv1a64:c31ecebc45bf27d4
% canonical-lessons: JA-C501-tsukimi, JA-C501-hanami, JA-C501-momiji, JA-C501-fuurin, JA-C501-chirasu

\chapter{Moon Viewing and Autumn Leaves}
\label{ch:ja-moon-viewing-and-autumn-leaves}
\hlchaptermodality{501}

\begin{chapteropening}
\textbf{By the end of this chapter:} \emph{I can talk about seasonal pastimes.}

Complete the last lesson of chapter 501: I can talk about seasonal pastimes.
\end{chapteropening}

\section[\texorpdfstring{\ja{つきみ} (tsukimi)}{つきみ (tsukimi)}]{\ja{つきみ} (tsukimi) — moon viewing}

\label{lesson:JA-C501-tsukimi}



\begin{quote}
Before the new one: say the Japanese for a riddle, then the Japanese for a slide (in a playground).
\end{quote}

\subsection*{You'll want to know: \ja{つきみ}}
\textbf{\ja{つきみ}} — \emph{tsukimi} — \textquotedblleft{}moon viewing\textquotedblright{}.

\begin{quote}
\textbf{\ja{あきの} \ja{ばん、つきみを} \ja{します。}} — \emph{aki no ban, tsukimi o shimasu} — \textquotedblleft{}On autumn evenings we go moon viewing.\textquotedblright{}
\end{quote}

\begin{cousinweb}
\textbf{\ja{つきみ}} (\emph{tsukimi}) is a native Japanese compound (wago), inherited rather than borrowed from Chinese: built from \emph{tsuki}, \textquotedblleft{}moon\textquotedblright{} and \emph{mi}, \textquotedblleft{}looking\textquotedblright{} (JMdict, KANJIDIC2).
\end{cousinweb}

\begin{grammarlens}[title={What you've built}]
One more word for describing what you lived and felt.
\end{grammarlens}

\subsection*{Your turn}
\begin{itemize}
\raggedright
  \item \emph{Say it:} \emph{tsukimi}
  \item \emph{Say it:} \emph{tsukimi}, once more
  \item \emph{Say it:} \emph{aki no ban, tsukimi o shimasu} — On autumn evenings we go moon viewing.
  \item \emph{Recall:} say the Japanese for a riddle, then the Japanese for a slide (in a playground), then say \emph{tsukimi} again
\end{itemize}

\begin{tcolorbox}[breakable,colback=teal!4,colframe=teal!35!black,title={Before you move on}]
What does \ja{つきみ} mean? (\textquotedblleft{}Moon viewing\textquotedblright{}.) Say it once more.
\end{tcolorbox}
\section[\texorpdfstring{\ja{はなみ} (hanami)}{はなみ (hanami)}]{\ja{はなみ} (hanami) — flower viewing; cherry-blossom viewing}

\label{lesson:JA-C501-hanami}



\begin{quote}
Before the new one: say the Japanese for a slide (in a playground), then the Japanese for moon viewing.
\end{quote}

\subsection*{You'll want to know: \ja{はなみ}}
\textbf{\ja{はなみ}} — \emph{hanami} — \textquotedblleft{}flower viewing; cherry-blossom viewing\textquotedblright{}.

\begin{quote}
\textbf{\ja{はるに} \ja{こうえんで} \ja{はなみを} \ja{します。}} — \emph{haru ni kōen de hanami o shimasu} — \textquotedblleft{}In spring we have a blossom-viewing party in the park.\textquotedblright{}
\end{quote}

\begin{cousinweb}
\textbf{\ja{はなみ}} (\emph{hanami}) is a native Japanese compound (wago), inherited rather than borrowed from Chinese: built from \emph{hana}, \textquotedblleft{}flower\textquotedblright{} and \emph{mi}, \textquotedblleft{}looking\textquotedblright{} (JMdict, KANJIDIC2).
\end{cousinweb}

\begin{grammarlens}[title={What you've built}]
One more word for describing what you lived and felt.
\end{grammarlens}

\subsection*{Your turn}
\begin{itemize}
\raggedright
  \item \emph{Say it:} \emph{hanami}
  \item \emph{Say it:} \emph{hanami}, once more
  \item \emph{Say it:} \emph{haru ni kōen de hanami o shimasu} — In spring we have a blossom-viewing party in the park.
  \item \emph{Recall:} say the Japanese for a slide (in a playground), then the Japanese for moon viewing, then say \emph{hanami} again
\end{itemize}

\begin{tcolorbox}[breakable,colback=teal!4,colframe=teal!35!black,title={Before you move on}]
What does \ja{はなみ} mean? (\textquotedblleft{}Flower viewing; cherry-blossom viewing\textquotedblright{}.) Say it once more.
\end{tcolorbox}
\section[\texorpdfstring{\ja{もみじ} (momiji)}{もみじ (momiji)}]{\ja{もみじ} (momiji) — autumn leaves; a maple}

\label{lesson:JA-C501-momiji}



\begin{quote}
Before the new one: say the Japanese for moon viewing, then the Japanese for flower viewing.
\end{quote}

\subsection*{You'll want to know: \ja{もみじ}}
\textbf{\ja{もみじ}} — \emph{momiji} — \textquotedblleft{}autumn leaves; a maple\textquotedblright{}.

\begin{quote}
\textbf{\ja{あきの} \ja{やまは} \ja{もみじが} \ja{きれいです。}} — \emph{aki no yama wa momiji ga kirei desu} — \textquotedblleft{}The autumn leaves on the mountains are beautiful.\textquotedblright{}
\end{quote}

\begin{cousinweb}
\textbf{\ja{もみじ}} (\emph{momiji}) is a native Japanese word (wago), inherited rather than borrowed from Chinese; its kanji was attached later as a kun'yomi, and its older origin is not settled (JMdict, KANJIDIC2).
\end{cousinweb}

\begin{grammarlens}[title={What you've built}]
One more word for describing what you lived and felt.
\end{grammarlens}

\subsection*{Your turn}
\begin{itemize}
\raggedright
  \item \emph{Say it:} \emph{momiji}
  \item \emph{Say it:} \emph{momiji}, once more
  \item \emph{Say it:} \emph{aki no yama wa momiji ga kirei desu} — The autumn leaves on the mountains are beautiful.
  \item \emph{Recall:} say the Japanese for moon viewing, then the Japanese for flower viewing, then say \emph{momiji} again
\end{itemize}

\begin{tcolorbox}[breakable,colback=teal!4,colframe=teal!35!black,title={Before you move on}]
What does \ja{もみじ} mean? (\textquotedblleft{}Autumn leaves; a maple\textquotedblright{}.) Say it once more.
\end{tcolorbox}
\section[\texorpdfstring{\ja{ふうりん} (fūrin)}{ふうりん (fūrin)}]{\ja{ふうりん} (fūrin) — a wind chime}

\label{lesson:JA-C501-fuurin}



\begin{quote}
Before the new one: say the Japanese for flower viewing, then the Japanese for autumn leaves.
\end{quote}

\subsection*{You'll want to know: \ja{ふうりん}}
\textbf{\ja{ふうりん}} — \emph{fūrin} — \textquotedblleft{}a wind chime\textquotedblright{}.

\begin{quote}
\textbf{\ja{なつに} \ja{ふうりんの} \ja{おとが} \ja{します。}} — \emph{natsu ni fūrin no oto ga shimasu} — \textquotedblleft{}In summer you hear the wind chimes.\textquotedblright{}
\end{quote}

\begin{cousinweb}
\textbf{\ja{ふうりん}} (\emph{fūrin}) is Sino-Japanese (kango): two Chinese characters read with their on'yomi, readings borrowed from Chinese. They are \emph{fū}, \textquotedblleft{}wind\textquotedblright{} (Middle Chinese \emph{pjuwng}) and \emph{rin}, \textquotedblleft{}bell\textquotedblright{} (\emph{leng}) (KANJIDIC2; Middle Chinese in Baxter's transcription).
\end{cousinweb}

\begin{grammarlens}[title={What you've built}]
One more word for describing what you lived and felt.
\end{grammarlens}

\subsection*{Your turn}
\begin{itemize}
\raggedright
  \item \emph{Say it:} \emph{fūrin}
  \item \emph{Say it:} \emph{fūrin}, once more
  \item \emph{Say it:} \emph{natsu ni fūrin no oto ga shimasu} — In summer you hear the wind chimes.
  \item \emph{Recall:} say the Japanese for flower viewing, then the Japanese for autumn leaves, then say \emph{fūrin} again
\end{itemize}

\begin{tcolorbox}[breakable,colback=teal!4,colframe=teal!35!black,title={Before you move on}]
What does \ja{ふうりん} mean? (\textquotedblleft{}A wind chime\textquotedblright{}.) Say it once more.
\end{tcolorbox}
\section[\texorpdfstring{\ja{ちらす} (chirasu)}{ちらす (chirasu)}]{\ja{ちらす} (chirasu) — to scatter (something); to strew}

\label{lesson:JA-C501-chirasu}



\begin{quote}
Before the new one: say the Japanese for autumn leaves, then the Japanese for a wind chime.
\end{quote}

\subsection*{You'll want to know: \ja{ちらす}}
\textbf{\ja{ちらす}} — \emph{chirasu} — \textquotedblleft{}to scatter (something); to strew\textquotedblright{}.

\begin{quote}
\textbf{\ja{かぜが} \ja{さくらの} \ja{はなを} \ja{ちらしました。}} — \emph{kaze ga sakura no hana o chirashimashita} — \textquotedblleft{}The wind scattered the cherry blossoms.\textquotedblright{}
\end{quote}

Polite \textbf{\ja{ちらします}}, \emph{chirashimasu}; a godan verb. Past (polite): \textbf{\ja{かぜが} \ja{はなを} \ja{ちらしました。}}, \emph{kaze ga hana o chirashimashita}.

\begin{cousinweb}
\textbf{\ja{ちらす}} (\emph{chirasu}) is a native Japanese word (wago), inherited rather than borrowed from Chinese, the transitive partner of \emph{chiru}, \textquotedblleft{}to fall (of leaves)\textquotedblright{}; its kanji was attached later as a kun'yomi (JMdict, KANJIDIC2).
\end{cousinweb}

\begin{grammarlens}[title={What you've built}]
One more word for describing what you lived and felt.
\end{grammarlens}

\subsection*{Your turn}
\begin{itemize}
\raggedright
  \item \emph{Say it:} \emph{chirasu}
  \item \emph{Say it:} \emph{chirasu}, once more
  \item \emph{Say it:} \emph{kaze ga sakura no hana o chirashimashita} — The wind scattered the cherry blossoms.
  \item \emph{Recall:} say the Japanese for autumn leaves, then the Japanese for a wind chime, then say \emph{chirasu} again
\end{itemize}

\begin{tcolorbox}[breakable,colback=teal!4,colframe=teal!35!black,title={Before you move on}]
What does \ja{ちらす} mean? (\textquotedblleft{}To scatter (something); to strew\textquotedblright{}.) Say it once more.
\end{tcolorbox}
